\section{Scope}

\subsection{Scope}
The intended users of the system are students, instructors, and administrators of a university. The core LMS functions, including course management, enrollment, assignments, quizzes, grades, and access control, are inherited from the open-source Moodle platform; this project does not rebuild these functions but instead focuses on extension and integration.
The project is carried out in two phases. In the specialized project phase, the team focuses on system analysis and design while implementing the basic features: data synchronization from Moodle, basic document processing, a minimal chatbot, academic early warning based on threshold rules or a simple model, and a chat interface together with a monitoring dashboard for instructors. In the graduation project phase, the team completes the system and improves the quality of the artificial intelligence components through hybrid retrieval, re-ranking, and an interpretable prediction model, while also optimizing performance and security and carrying out testing, evaluation, and deployment.
With respect to artificial intelligence features, the scope of the project covers a chatbot that answers questions based on course materials and an early-warning function for students at risk of poor academic performance. A course recommendation feature for students is planned for the second phase, subject to the available schedule.

\subsection{Limitations}
Because the Moodle core is not modified and only plugins and external services are developed, the features of the system depend on the data and programming interfaces that Moodle provides. Regarding data, since the team does not have access to real learning data from the university, the warning model is built on public datasets or simulated data; consequently, the evaluation results are experimental and do not fully reflect real-world conditions.
Regarding documents and language, the chatbot supports common formats such as PDF, DOCX, and PPTX, while low-quality scanned documents, complex tables, and mathematical formulas are handled only to a limited extent. The system prioritizes Vietnamese, and its quality in other languages lies outside the scope of evaluation. Regarding resources, the language model runs on limited hardware and budget, so latency and answer quality may be lower than those of large-scale commercial services. The system is also tested only at a small scale, with a few courses and sample users, and has not been validated at the scale of an entire university.
Furthermore, features such as camera-based cheating detection, a dedicated mobile application, and integration with the university's official credit registration system are beyond the scope of this project. Finally, it should be noted that the outputs of the chatbot and the warning model are intended solely as supporting references and do not replace the decisions of instructors.
